% ============================================================
% ExaBench v0.1 — arXiv preprint
% ============================================================
\documentclass[12pt]{article}

\usepackage[margin=1in]{geometry}
\usepackage{amsmath}
\usepackage{amssymb}
\usepackage{booktabs}
\usepackage{microtype}
\usepackage[colorlinks=true,linkcolor=blue,citecolor=blue,urlcolor=blue]{hyperref}
\usepackage[capitalize,nameinlink]{cleveref}
\usepackage[numbers,sort&compress]{natbib}

\title{\textbf{ExaBench: A Role-Aware, RBAC-Enforced Benchmark\\
       for LLM Agents in High-Performance Computing Operations}}

\author{%
  Anonymous Authors\\[2pt]
  \small Submitted for review
}

\date{March 2026}

\begin{document}

\maketitle

% -----------------------------------------------------------------------
\begin{abstract}
LLM agents are increasingly deployed in operational domains, yet no benchmark
exists for high-performance computing (HPC) facilities---environments governed
by strict role hierarchies, scheduler-native tooling, and legally binding
access-control policies where a single privilege violation can cascade into
cluster-wide disruption.
We introduce \textbf{ExaBench}, the first role-aware, RBAC-enforced benchmark
for HPC agent evaluation.
ExaBench provides 30 tasks across a $3{\times}3$ role--QCAT grid
(\textsc{User}/\textsc{Operator}/\textsc{Manager} $\times$
\textsc{Job}/\textsc{Mon}/\textsc{Energy}) backed by a deterministic mock HPC
environment of five tools and 20 environment snapshots.
Governance is evaluated via a hard-fail scorer and the
Completion-under-Policy (CuP) metric; deployment readiness is reported on the
five-dimensional CLEAR scorecard.
Evaluating GPT-4o against a zero-tool baseline, we find GPT-4o achieves a
0.517 aggregate score (+53\% over baseline) but zero CLEAR Assurance,
indicating complete RBAC compliance failure on every tool-using episode.
Reliability experiments further reveal a binary competency pattern: hard tasks
are solved with $\text{pass}^8=1.000$ while easy and medium tasks fail
consistently, implicating RBAC-aware reasoning rather than raw capability as
the critical gap.
Benchmark, environments, and evaluation code are open-sourced.
\end{abstract}

% -----------------------------------------------------------------------
\section{Introduction}
\label{sec:introduction}

High-performance computing (HPC) facilities are among the most operationally
complex environments that software agents could be deployed in.
A modern supercomputing centre presents a tightly coupled system of
scheduler daemons, telemetry pipelines, energy controllers, and policy
enforcement layers, all governed by a strict role hierarchy in which the
same action---cancelling a job, reading power telemetry, modifying a
queue---may be a routine duty for an operator and an impermissible
privilege escalation for a scientific user.
Mistakes in this environment are not merely costly; they can cascade into
cluster-wide disruption, security incidents, or violations of facility
governance obligations that are legally binding.

Despite this operational reality, the rapid deployment of LLM-based agents
across technical domains has not been accompanied by a commensurate evaluation
infrastructure for the HPC context.
Existing agent benchmarks have been designed for software engineering
(SWE-bench), factual web tasks (GAIA), desktop GUI interaction (OSWorld),
and cloud operations (Cloud-OpsBench, $\tau$-bench).
While these works have collectively established the methodological vocabulary
of agent evaluation---deterministic state injection, execution-verified
ground truth, policy-following paradigms---none exposes HPC-native tooling
(SLURM job management, Prometheus-style telemetry, facility energy APIs),
none enforces role-contingent correctness in which different agent personas
must produce different valid answers from the same environment snapshot,
and none treats privilege violation as an automatic, non-negotiable scoring
failure rather than a soft penalty.
No existing benchmark can answer the question: \emph{is this LLM agent safe
and capable enough to operate in a production HPC facility?}

We introduce \textbf{ExaBench}, the first role-aware, RBAC-enforced benchmark
for evaluating LLM agents in HPC operations.
ExaBench couples a curated 30-task corpus spanning three user roles, three
query categories (QCATs), and three difficulty tiers with a deterministic
mock HPC environment, a governance hard-fail scoring model, and the CLEAR
five-dimension evaluation framework.
RBAC compliance is elevated from a soft rubric dimension to a binary
outcome gate via the Completion-under-Policy (CuP) metric, ensuring that
a task answered correctly through an unauthorised tool call receives a
score of zero.

\paragraph{Contributions.}
\begin{itemize}
  \item \textbf{First role-aware HPC agent benchmark.} ExaBench provides
    30 tasks across a $3 \times 3$ role--QCAT grid
    (\textsc{User}/\textsc{Operator}/\textsc{Manager}
    $\times$ \textsc{Job}/\textsc{Mon}/\textsc{Energy}) at three difficulty
    tiers (Easy / Medium / Hard), with role-differentiated ground truth
    derived from deterministic environment snapshots.
  \item \textbf{RBAC hard-fail scoring with CuP.} Any privilege violation
    in an agent's execution trace zeros the task score regardless of
    outcome quality, instantiated via the Completion-under-Policy (CuP)
    metric adapted from ST-WebAgentBench~\cite{st-webagentbench}.
  \item \textbf{CLEAR evaluation framework.} All agents are assessed across
    five dimensions---Cost, Latency, Efficacy, Assurance, and
    Reliability---providing a deployment-readiness signal that correlates
    more strongly with production suitability than accuracy alone
    \citep{clear}.
  \item \textbf{Finding: RBAC is the critical gap.}  GPT-4o achieves a
    0.517 aggregate score and strong tool-use performance (0.905) on the
    21-task dev split, yet its CLEAR Assurance score is 0.000, indicating
    zero-tolerance RBAC compliance failure across all tool-using episodes.
  \item \textbf{Finding: binary competency across difficulty tiers.}
    Reliability experiments reveal that GPT-4o either solves tasks
    deterministically ($\text{pass}^8 = 1.000$) or fails deterministically
    ($\text{pass}^8 = 0.000$), with hard tasks achieving perfect reliability
    while easy and medium tasks fail consistently---a difficulty-inverted
    pattern that implicates RBAC-aware reasoning rather than raw capability.
\end{itemize}

The remainder of the paper is organised as follows.
\Cref{sec:related-work} surveys related benchmarks and prior HPC AI work.
\Cref{sec:benchmark} describes the ExaBench task corpus, mock environment,
governance model, and scoring framework.
\Cref{sec:evaluation} reports experimental results.
\Cref{sec:discussion} analyses key findings and implications for deployment.
\Cref{sec:conclusion} summarises contributions and future directions.

% -----------------------------------------------------------------------
\section{Related Work}
\label{sec:related-work}

The rapid proliferation of LLM-based agents has driven a parallel expansion
of agent evaluation frameworks, from single-turn question-answering to
multi-step, tool-using episodes in interactive
environments~\citep{llm-agent-eval-survey}.
Existing benchmarks, however, have largely been designed for web, software,
or enterprise domains, leaving the unique demands of high-performance
computing (HPC) operations---multi-tier role hierarchies, scheduler-native
tooling, and policy-gated resource access---systematically unaddressed.
ExaBench is designed to fill this gap by unifying role-aware task design,
RBAC hard-fail evaluation, and HPC-native tooling under a single,
deterministic benchmark framework.

\subsection*{General Agent Benchmarks}

SWE-bench~\citep{swe-bench} operationalised software engineering as a
patch-generation task over frozen GitHub repositories, establishing the
pattern of execution-verified ground truth.
GAIA~\citep{gaia} extended this to factual multi-step reasoning with tool
use, while OSWorld~\citep{osworld} moved into desktop GUI environments with
VM-based deterministic reset.
AgentBench~\citep{agentbench} aggregated eight task environments---including
OS shells and web browsers---into a unified evaluation harness.
TheAgentCompany~\citep{theagentcompany} introduced a simulated corporate
intranet (GitLab, Plane, RocketChat) with role labels such as \emph{Admin}
and \emph{SDE}, representing the closest prior work to role-differentiated
evaluation.
Despite these advances, none of these benchmarks expose SLURM job
management, HPC telemetry, or energy-monitoring tooling; none enforce
role-contingent correctness (where the same action is correct for a
\texttt{sysadmin} but forbidden for a \texttt{scientific\_user}); and none
treat privilege violation as an automatic hard-fail.
ExaBench replaces generic environments with an HPC-native mock stack and
elevates access-control compliance from a soft rubric dimension to a binary,
non-negotiable outcome gate.

\subsection*{Operations and Cloud Benchmarks}

Cloud-OpsBench~\citep{cloud-opsbench} introduced the snapshot-driven
evaluation pattern for cloud operations, injecting deterministic Kubernetes
and cloud-API state before each episode.
$\tau$-bench~\citep{tau-bench} contributed the
\emph{pass\textsuperscript{k}} reliability metric---measuring the
probability that an agent succeeds consistently across $k$ independent
runs---and a policy-following paradigm in which agents must comply with
domain-specific rules encoded in a natural-language document.
ST-WebAgentBench~\citep{st-webagentbench} introduced the
Completion-under-Policy (CuP) scorer, which zeroes out task scores
on any policy violation across six constraint dimensions.
These benchmarks collectively establish the methodological vocabulary that
ExaBench adopts and extends: deterministic state injection from
$\tau$-bench and Cloud-OpsBench, hard-fail policy enforcement from
ST-WebAgentBench, and $\text{pass}^k$ reliability from $\tau$-bench.
However, all three operate exclusively in cloud or retail domains, lack
HPC-native tools (SLURM, Prometheus, facility DCIM), provide no multi-role
task variants in which different roles receive different valid answers from
the same environment snapshot, and do not model the three-tier HPC role
hierarchy (\texttt{scientific\_user} $<$ \texttt{sysadmin} $<$
\texttt{facility\_admin}) that governs real facility operations.

\subsection*{HPC-Specific AI Work}

HPC-GPT~\citep{hpc-gpt} fine-tuned a language model for HPC-domain question
answering over PLP and MLPerf corpora and for static data-race detection,
demonstrating that HPC-specific pretraining improves domain accuracy.
ChatHPC~\citep{chathpc} deployed a RAG-backed chatbot pipeline on Frontier
to answer HPC usage questions and generate job scripts.
HPCAgentTester~\citep{hpcagenttester} proposed a multi-agent recipe for
generating parallel-program test cases, targeting race conditions and
deadlocks in MPI and OpenMP code.
These works establish that HPC is a meaningful LLM target domain, but none
constitutes an \emph{agent evaluation benchmark}: HPC-GPT and ChatHPC are
chatbot systems with no tool-calling, no ground-truth scoring, and no
multi-role evaluation; HPCAgentTester generates tests rather than assessing
operational agents.
None covers the MON, ENERGY, SEC, or FAC query categories (QCATs) that
characterise real facility operations, and none measures agent reliability
across repeated runs or enforces RBAC compliance as a scoring dimension.
ExaBench addresses each of these gaps by providing a structured task corpus
across ten HPC QCATs, a CLEAR five-dimension scorecard~\citep{clear}
including pass\textsuperscript{k} reliability and RBAC assurance, and an
HPC-specific error taxonomy derived from the TRAIL framework~\citep{trail}.

\subsection*{Feature Comparison}

\Cref{tab:related-work} summarises the key differentiating features of
ExaBench relative to prior benchmarks.

\begin{table}[htbp]
\centering
\caption{Feature comparison of agent benchmarks.
  \textbf{Role-aware}: tasks require different correct answers per agent
  role.
  \textbf{RBAC hard-fail}: any privilege violation automatically scores
  zero.
  \textbf{HPC tools}: benchmark exposes SLURM, telemetry, or
  energy-monitoring APIs.
  \textbf{CLEAR metric}: evaluation reports all five CLEAR dimensions
  (Cost, Latency, Efficacy, Assurance, Reliability).
  \textbf{Det.\ snapshot}: environment state is deterministically
  reproduced per episode.}
\label{tab:related-work}
\small
\setlength{\tabcolsep}{5pt}
\begin{tabular}{lccccc}
\toprule
\textbf{Benchmark}
  & \textbf{Role-aware}
  & \textbf{RBAC hard-fail}
  & \textbf{HPC tools}
  & \textbf{CLEAR metric}
  & \textbf{Det.\ snapshot} \\
\midrule
SWE-bench~\citep{swe-bench}
  & $\times$ & $\times$ & $\times$ & $\times$ & $\times$ \\
GAIA~\citep{gaia}
  & $\times$ & $\times$ & $\times$ & $\times$ & $\times$ \\
OSWorld~\citep{osworld}
  & $\times$ & $\times$ & $\times$ & $\times$ & $\circ$  \\
TheAgentCompany~\citep{theagentcompany}
  & $\times$ & $\times$ & $\times$ & $\times$ & $\checkmark$ \\
\midrule
Cloud-OpsBench~\citep{cloud-opsbench}
  & $\times$ & $\times$ & $\times$ & $\times$ & $\checkmark$ \\
$\tau$-bench~\citep{tau-bench}
  & $\times$ & $\times$ & $\times$ & $\circ$  & $\times$ \\
ST-WebAgentBench~\citep{st-webagentbench}
  & $\times$ & $\circ$  & $\times$ & $\times$ & $\times$ \\
\midrule
HPC-GPT~\citep{hpc-gpt}
  & $\times$ & $\times$ & $\times$ & $\times$ & $\times$ \\
ChatHPC~\citep{chathpc}
  & $\times$ & $\times$ & $\circ$  & $\times$ & $\times$ \\
HPCAgentTester~\citep{hpcagenttester}
  & $\times$ & $\times$ & $\circ$  & $\times$ & $\times$ \\
\midrule
\textbf{ExaBench (ours)}
  & $\checkmark$ & $\checkmark$ & $\checkmark$ & $\checkmark$ & $\checkmark$ \\
\bottomrule
\multicolumn{6}{l}{$\checkmark$\,=\,fully supported;\;
                   $\circ$\,=\,partially supported;\;
                   $\times$\,=\,absent.}
\end{tabular}
\end{table}

% -----------------------------------------------------------------------
\section{Benchmark Design}
\label{sec:benchmark}

\subsection{Task Design}
\label{sec:task-design}

ExaBench comprises 30 tasks organised along two primary axes: \emph{user role} and
\emph{query category} (QCAT).  Three roles capture the principal stakeholder
personas of an HPC facility: \textsc{User} (a standard scientific user whose
visibility is restricted to their own jobs), \textsc{Operator} (a system
administrator with facility-wide read and write access), and \textsc{Manager}
(a facility director with full access including energy budget and policy data).
Three QCATs reflect the operational concerns most frequently surfaced in HPC
monitoring and management literature: \textsc{Job} (SLURM job lifecycle:
submission, queuing, execution, failure), \textsc{Mon} (node telemetry:
CPU/GPU/memory utilisation, anomaly detection, node health), and
\textsc{Energy} (power draw, energy consumption, efficiency, budget tracking).
The role--QCAT pairing forms a $3 \times 3$ grid (\Cref{fig:taxonomy}), with
each cell populated by tasks at three difficulty tiers: \emph{Easy}
(single-tool, point-lookup, OLTP-style), \emph{Medium} (multi-step sequential
tool chaining, OLAP-style aggregation), and \emph{Hard} (parallel tool calls,
cross-QCAT reasoning, or rubric-scored open-ended diagnosis).  Difficulty is
defined operationally following the Berkeley Function-Calling Leaderboard
(BFCL) taxonomy~\citep{bfcl2024}: easy tasks require one tool call with
explicit arguments; medium tasks require chained calls where the output of one
call provides arguments for the next; hard tasks require parallel invocations
or free-form expert-level recommendations scored by an LLM judge.

Each task is specified as a \texttt{TaskSpec} object carrying a
natural-language prompt, the expected role, a list of \texttt{allowed\_tools},
optional \texttt{expected\_tool\_calls} for deterministic scoring, and a ground
truth answer derived from the associated environment snapshot.  Tasks with
verifiable numeric or categorical answers use a \emph{deterministic} scoring
path; diagnostic and advisory tasks use a \emph{rubric} scoring path evaluated
by an LLM judge against a structured criterion set.

\begin{figure}[t]
  \centering
  \fbox{\parbox{0.85\linewidth}{\centering
    \textbf{Figure 1:} ExaBench task taxonomy --- $3 \times 3$ Role $\times$
    QCAT grid (User / Operator / Manager vs.\ Job / Mon / Energy) with
    difficulty tiers (Easy / Medium / Hard) shown per cell.\\[4pt]
    \emph{[Figure to be inserted: heat-map style grid, 3 rows $\times$ 3 cols,
    each cell labelled with task count per tier.]}
  }}
  \caption{ExaBench task taxonomy: the $3 \times 3$ role--QCAT grid.  Each
  cell contains tasks at three difficulty tiers.  Role-gated QCATs (e.g.,
  Energy tasks visible only to Manager) are indicated by shaded borders.}
  \label{fig:taxonomy}
\end{figure}

\subsection{Mock HPC Environment}
\label{sec:mock-env}

Each task episode is grounded in a \emph{deterministic mock HPC environment}
consisting of five mock tools and an environment snapshot.  The five tools are:
\texttt{slurm} (job and partition queries), \texttt{telemetry} (time-series
node metrics), \texttt{rbac} (role and permission lookups), \texttt{facility}
(energy and power data), and \texttt{docs} (policy document retrieval).  Tool
signatures follow the OpenAI function-calling format~\citep{bfcl2024} and are
defined in a canonical catalog (\texttt{hpc\_tool\_catalog.yaml}) that
specifies required and optional parameters, valid enumerations, return shapes,
and per-role visibility, following the \texttt{required\_parameters}\,/\,\texttt{optional\_parameters}
split introduced by ToolBench~\citep{qin2023toolllm}.

The benchmark ships 20 environment snapshots---static JSON bundles capturing
a consistent point-in-time state of a simulated HPC cluster including job
queue state, node health, telemetry time series, energy readings, and RBAC
policies.  Snapshots are versioned and immutable.  All tool calls are
\emph{deterministic}: given the same snapshot and the same tool invocation,
the returned observation is identical across runs, evaluators, and time.
Determinism is a deliberate design choice motivated by two requirements.
First, it decouples evaluation correctness from the stochasticity of
real-cluster state, enabling fair model-to-model comparison.  Second, it
enables reproducible pass\textsuperscript{$k$} reliability estimation
(\Cref{sec:scoring}): the variance in repeated runs reflects the
LLM's sampling stochasticity, not environmental noise.  Agents terminate
each episode by invoking a mandatory \texttt{finish} tool, following the
ToolBench design~\citep{qin2023toolllm}; episodes without a \texttt{finish}
call are recorded as agent abandonment and receive an outcome score of zero.

\subsection{Governance and RBAC}
\label{sec:governance}

Role-based access control is a first-class evaluation dimension in ExaBench,
not a post-hoc annotation.  Every environment bundle contains a
machine-readable \texttt{rbac\_policy.yaml} (schema version 1.1) that specifies,
per role, the list of \texttt{allowed\_tools}, per-partition walltime limits,
and data-access tiers (tier~1 public through tier~4 sensitive).  An
agent-readable \texttt{rbac\_policy.md} document in the environment's
\texttt{docs/} directory mirrors this policy in natural language, following the
$\tau$-bench paradigm~\citep{tau-bench} in which the policy document is a
separate artifact the agent must consult rather than a fact baked into the
system prompt.

Governance evaluation distinguishes three violation severity levels.  A
\emph{hard fail} occurs when the execution trace carries a
\texttt{hard\_fail\_reason} containing a permission-related cause; the
governance score is forced to 0.0 and the episode is terminated.  A
\emph{forbidden tool call} occurs when the agent invokes a tool not in the
task's \texttt{allowed\_tools} set (penalty: $-0.50$ per call), following the
BFCL forbidden-call detection paradigm~\citep{bfcl2024}.  A \emph{soft
violation} occurs when a tool call returns a \texttt{permission\_denied}
observation (penalty: $-0.25$ per occurrence).  The \texttt{GovernanceScorer}
aggregates these into a graded score $g \in [0,1]$.

Beyond graded governance scoring, ExaBench reports the
\textbf{Completion-under-Policy (CuP)} metric, adapted from
ST-WebAgentBench~\citep{st-webagentbench}.  CuP applies a binary gate to
task-level outcome scores:
\begin{equation}
  \mathrm{CuP}_t \;=\; C_t \times \mathbf{1}\!\left[\textstyle\sum_d V^d_t = 0\right]
  \label{eq:cup}
\end{equation}
where $C_t \in [0,1]$ is the raw completion score and $V^d_t$ is a violation
indicator across six policy dimensions.  A task that is answered correctly but
involves any policy violation receives a CuP score of zero, regardless of the
magnitude of the violation.  The gap between raw completion rate and aggregate
CuP is a direct measure of safety-completeness failure.

\subsection{Scoring Framework}
\label{sec:scoring}

ExaBench adopts the \textbf{CLEAR} framework~\citep{clear} as its top-level
evaluation structure, extending task-outcome accuracy with four additional
dimensions that correlate more strongly with production deployment readiness
($\rho = 0.83$ vs.\ $\rho = 0.41$ for accuracy alone).  The five CLEAR
dimensions as instantiated in ExaBench are:

\begin{itemize}
  \item \textbf{Efficacy (E):} task outcome quality (0--1), computed by the
    \texttt{HybridScorer} combining deterministic ground-truth matching and
    LLM-judge rubric scoring; CuP-gated for governance-violating episodes.
  \item \textbf{Assurance (A):} fraction of tasks with no RBAC violations
    ($g = 1.0$), reported as the binary compliance rate across the task set.
  \item \textbf{Reliability (R):} pass\textsuperscript{$k$} consistency
    estimated via the $\tau$-bench unbiased estimator over $n$ repeated runs.
  \item \textbf{Cost (C):} per-task API spend in USD (tokens $\times$ model
    price), min-max normalised across compared models at report time.
  \item \textbf{Latency (L):} wall-clock time from task start to
    \texttt{finish} call, min-max normalised (lower is better).
\end{itemize}

\noindent The composite CLEAR score is:

\begin{equation}
  \mathrm{CLEAR} \;=\; 0.2\,C_{\mathrm{norm}}
                     + 0.2\,L_{\mathrm{norm}}
                     + 0.2\,E
                     + 0.2\,A
                     + 0.2\,R
  \label{eq:clear}
\end{equation}

\noindent where $C_{\mathrm{norm}}$ and $L_{\mathrm{norm}}$ are min-max
normalised so that lower cost and lower latency map to higher scores, and
$E, A, R \in [0,1]$ require no further normalisation.  Equal weights (0.2
each) follow the original CLEAR paper defaults; domain-specific weight
overrides are supported via named scoring profiles.  All five dimension scores
are stored per task in \texttt{BenchmarkResult} and aggregated at the run
level to produce the CLEAR scorecard reported on the leaderboard.

% -----------------------------------------------------------------------
\section{Evaluation}
\label{sec:evaluation}

\subsection{Experimental Setup}

We evaluate two agents on the 21-task ExaBench dev split: \textbf{GPT-4o}
(Azure OpenAI) and a \textbf{direct\_qa} zero-tool baseline that answers from
parametric knowledge without invoking any tools. All tool calls are handled by
mock HPC implementations (\texttt{slurm\_tool}, \texttt{telemetry\_tool},
\texttt{energy\_tool}) backed by fixture data; no real cluster is required.
The 9-task test split is held out and not evaluated in this study. Agents are
scored on five dimensions---Outcome, Tool Use, Governance, Efficiency, and
Grounding---combined into a weighted Aggregate score.

\subsection{Main Results}

\Cref{tab:main} reports mean scores across all 21 dev tasks. GPT-4o
outperforms direct\_qa on Aggregate (0.517 vs.\ 0.337), driven by strong
Tool Use (0.905) and Efficiency (0.949). The Governance dimension reverses
this ordering sharply (0.095 vs.\ 1.000), a phenomenon we examine in
\Cref{sec:discussion}.

\begin{table}[ht]
\centering
\caption{Mean scores across the 21-task dev split. Higher is better for all dimensions.}
\label{tab:main}
\begin{tabular}{lcccccc}
\toprule
\textbf{Model} & \textbf{Aggregate} & \textbf{Outcome} & \textbf{Tool Use}
  & \textbf{Governance} & \textbf{Efficiency} & \textbf{Grounding} \\
\midrule
GPT-4o     & 0.517 & 0.448 & 0.905 & 0.095 & 0.949 & 0.561 \\
direct\_qa & 0.337 & 0.248 & 0.000 & 1.000 & 1.000 & 0.000 \\
\bottomrule
\end{tabular}
\end{table}

\subsection{CLEAR Scorecard}

\Cref{tab:clear} presents the CLEAR composite (Cost, Latency, Efficacy,
Assurance, Reliability) for both agents. GPT-4o achieves CLEAR~=~0.324;
positive Efficacy ($E=0.448$) and Reliability ($R=0.619$) are undermined by
zero Assurance ($A=0.000$), which captures RBAC compliance alongside outcome
quality. The direct\_qa baseline is not assigned a CLEAR composite because it
incurs no token cost, latency, or tool-use risk, making the cost-normalised
components undefined; its $E$ and $A$ values are shown for reference.

\begin{table}[ht]
\centering
\caption{CLEAR scorecard. CNA = cost-normalised accuracy (outcome / cost
  $\times$ 100); CPS = cost per session in USD. Dashes indicate undefined
  quantities for the zero-tool baseline.}
\label{tab:clear}
\begin{tabular}{lcccccc}
\toprule
\textbf{Model} & \textbf{CLEAR} & \textbf{E} & \textbf{A} & \textbf{R}
  & \textbf{CNA} & \textbf{CPS (\$)} \\
\midrule
GPT-4o     & 0.324 & 0.448 & 0.000 & 0.619 & 6536.6 & \$0.0159 \\
direct\_qa & ---   & 0.248 & 1.000 & 0.000 & ---    & ---      \\
\bottomrule
\end{tabular}
\end{table}

\subsection{Role $\times$ QCAT Stratification}

\Cref{tab:role_qcat} breaks GPT-4o scores by the three user roles and three
task categories in the dev split. Sysadmin tasks yield the highest scores
across all QCATs (0.503--0.625), while scientific\_user tasks are consistently
lower (0.383--0.435). Cell sizes are small ($n=1$--5), so these figures are
indicative rather than statistically conclusive.

\begin{table}[ht]
\centering
\caption{GPT-4o mean aggregate score by role $\times$ QCAT.
  $n$ denotes the number of tasks in each cell.}
\label{tab:role_qcat}
\begin{tabular}{lccc}
\toprule
\textbf{Role} & \textbf{ENERGY} & \textbf{JOB} & \textbf{MON} \\
\midrule
scientific\_user & 0.383 ($n=1$) & 0.435 ($n=3$) & 0.416 ($n=1$) \\
sysadmin         & 0.625 ($n=1$) & 0.614 ($n=2$) & 0.503 ($n=4$) \\
facility\_admin  & 0.586 ($n=5$) & 0.391 ($n=2$) & 0.584 ($n=2$) \\
\bottomrule
\end{tabular}
\end{table}

\subsection{Reliability}

\Cref{tab:reliability} shows $\text{pass}^8$ scores from $n=10$ repeated runs
on five representative tasks, sampled to cover all three QCATs and three
difficulty levels. The $\text{pass}^8$ estimator~\citep{tau-bench} treats a
run as passing when aggregate score exceeds 0.5. The spread across tasks
(0.000--1.000) satisfies validity gate~V4.

\begin{table}[ht]
\centering
\caption{GPT-4o reliability on five tasks ($n=10$ runs each, pass
  threshold~=~0.5). $\text{pass}^8$ uses the \citet{tau-bench} unbiased
  estimator; Mean and Std are per-run aggregate scores.}
\label{tab:reliability}
\begin{tabular}{llllllll}
\toprule
\textbf{Task} & \textbf{QCAT} & \textbf{Difficulty}
  & \textbf{$n_\text{runs}$} & \textbf{$n_\text{passed}$}
  & \textbf{pass$^8$} & \textbf{Mean} & \textbf{Std} \\
\midrule
JOB\_USR\_001    & JOB    & easy   & 10 & 0  & 0.000 & 0.422 & 0.001 \\
JOB\_SYS\_002    & JOB    & medium & 10 & 1  & 0.000 & 0.440 & 0.041 \\
MON\_SYS\_003    & MON    & medium & 10 & 0  & 0.000 & 0.373 & 0.005 \\
ENERGY\_FAC\_003 & ENERGY & hard   & 10 & 10 & 1.000 & 0.716 & 0.007 \\
MON\_SYS\_006    & MON    & hard   & 10 & 10 & 1.000 & 0.587 & 0.005 \\
\bottomrule
\end{tabular}
\end{table}

% -----------------------------------------------------------------------
\section{Discussion}
\label{sec:discussion}

\subsection{The RBAC Paradox}

The most striking result is the inversion of Governance scores: GPT-4o scores
0.095 while direct\_qa scores 1.000. This is not a failure of GPT-4o but a
structural property of the benchmark. A zero-tool agent cannot violate RBAC
because it never calls tools; its perfect Governance score is vacuous rather
than a meaningful safety signal. Conversely, GPT-4o's active tool use exposes
it to policy constraints that the baseline never encounters. The CLEAR
Assurance dimension ($A=0.000$ for GPT-4o, $A=1.000$ for direct\_qa) makes
this trade-off legible: a model with high Outcome but zero Assurance is not
deployable in a real HPC environment where privilege boundaries are legally
and operationally binding. ExaBench surfaces this tension by design---the
benchmark should not be ``gamed'' by avoiding tool use.

\subsection{Binary Competency}

The reliability results reveal a clean capability threshold with no
intermediate regime. GPT-4o either solves tasks deterministically
(ENERGY\_FAC\_003 and MON\_SYS\_006: $\text{pass}^8=1.000$, 10/10 runs
passing, std $\leq 0.007$) or fails deterministically (JOB\_USR\_001,
MON\_SYS\_003: $\text{pass}^8=0.000$, 0/10 passing, std $\leq 0.005$). The
single partial pass in JOB\_SYS\_002 ($n_\text{passed}=1$,
$\text{pass}^8=0.000$) is consistent with the estimator's strict binomial
penalty rather than genuine stochasticity. Counterintuitively, this pattern
is difficulty-inverted: the two hard tasks show perfect reliability while
easy and medium tasks fail reliably. We hypothesise that hard tasks in
ExaBench v0.1 involve structured data-retrieval workflows well-represented in
GPT-4o's training, while easy tasks require nuanced RBAC-aware reasoning that
the model consistently under-applies.

\subsection{Role Stratification}

GPT-4o gains over direct\_qa are strongly stratified by role. Sysadmins
benefit most, with deltas of $+0.28$--$+0.30$ across ENERGY and JOB
categories, consistent with the hypothesis that operator-level tasks (queue
management, node monitoring) align well with the procedural tool-use patterns
GPT-4o has learned. Scientific users show the smallest gains
($+0.04$--$+0.10$): their tasks require interpreting job output in domain
context, a grounding capability where the model's advantage over parametric
knowledge is limited. Facility administrators occupy an intermediate position
($+0.06$--$+0.25$). These patterns suggest that benchmark coverage weighted
toward operator roles would overstate the practical utility of current LLMs
for the scientific user population, which is the largest user class on
production HPC systems.

\subsection{CLEAR as a Deployment Readiness Signal}

A CLEAR score of 0.324 for GPT-4o should be interpreted as a signal that the
model is not yet deployable in a production HPC setting, despite positive
Efficacy and Reliability components. The bottleneck is Assurance ($A=0.000$):
the CuP scoring model maps any model exhibiting non-compliant tool-use
patterns to zero Assurance, regardless of outcome performance. This reflects
an operational reality in HPC environments---a single privilege escalation or
unauthorised queue modification can cascade into system-wide disruption.
Future work should focus on interventions that improve Governance without
sacrificing the Tool Use gains that drive Outcome performance.

% -----------------------------------------------------------------------
\section{Limitations}
\label{sec:limitations}

\paragraph{Partial QCAT coverage.}
Only three of the ten planned task categories (JOB, MON, ENERGY) are included
in v0.1. The remaining seven categories (PERF, DATA, SEC, ARCH, AIOPS, DOCS,
and one reserved set) are not yet populated. Results may not transfer to
tasks involving security auditing, data management, or architectural
configuration.

\paragraph{Single model evaluated.}
Only Azure GPT-4o was evaluated in this study. Claude Sonnet~4.6 and
GPT-4o-mini were originally planned but excluded due to API access
constraints. The RBAC paradox and binary competency patterns described above
may be specific to GPT-4o's training distribution; multi-model comparison is
deferred to v0.2.

\paragraph{Mock tool fidelity.}
All tool calls are served by fixture-backed mock implementations. Real SLURM
schedulers exhibit failure modes---partial job arrays, preemption races,
filesystem quota errors---that deterministic fixtures cannot capture. Agents
evaluated on ExaBench v0.1 may over- or under-perform relative to a live
cluster depending on sensitivity to these failure modes.

\paragraph{Test split not evaluated.}
The 9-task test split is held out and has not been evaluated by any model.
Reported scores reflect dev-set performance only; test-split evaluation is
planned for v0.2 alongside additional model entries.

\paragraph{Small role~$\times$~QCAT cell sizes.}
Several cells in \Cref{tab:role_qcat} contain only a single task ($n=1$).
Differences between roles within a QCAT should be treated as preliminary
observations rather than statistically reliable findings.

% -----------------------------------------------------------------------
\section{Conclusion}
\label{sec:conclusion}

ExaBench is the first benchmark to jointly address role-aware task design,
RBAC hard-fail evaluation, and HPC-native tooling for LLM agent assessment.
Our v0.1 evaluation establishes two key findings: active tool use is
necessary for high Outcome scores yet currently incompatible with RBAC
compliance, and agent competency is binary across difficulty tiers rather
than graded. The CLEAR framework makes both findings legible as a
deployment-readiness signal. ExaBench v0.2 will expand to all ten planned
QCATs (adding PERF, DATA, SEC, ARCH, AIOPS, and DOCS), introduce Claude
Sonnet and open-weight models (Llama-3, Mistral), evaluate the held-out test
split, and launch a public leaderboard tracking five-dimensional CLEAR scores
across the growing model landscape.

% -----------------------------------------------------------------------
\bibliographystyle{plainnat}
\bibliography{references}

\end{document}
